% pdflatex-tate-math-rawbox.tex -- P6: raw primitive box 境界の保証
% 判定（TATE-PASS/TATE-FAIL; runner は -halt-on-error なしで実行）:
%   LEGAL   : \text{...} 引数内の raw \hbox は合法（wrapper exempt）で
%             rawerr カウンタが立たないこと
%   RAWERR  : 捕捉本体トップレベルの raw \hbox で ClassError 発火
%             （\jlreqpdf@math@rawerrcnt で実測; 文言は runner が grep）
%   VBOXERR : raw \vbox でも同様（カウンタ合計 2）
\documentclass[pdflatex,tate,paper=a5]{jlreq}
\usepackage{amsmath}
\pagestyle{empty}
\makeatletter
\begin{document}
% ---- LEGAL: wrapper 引数内の raw box は検査対象外 ----
漢$\text{和\hbox{文}}$漢
\ifnum\jlreqpdf@math@rawerrcnt=0
  \typeout{TATE-PASS: LEGAL (raw box in wrapper arg: no ClassError)}%
\else
  \typeout{TATE-FAIL: LEGAL (unexpected raw err count=\the\jlreqpdf@math@rawerrcnt)}%
\fi
% ---- RAWERR: 捕捉本体のトップレベル raw \hbox → ClassError ----
漢$x\hbox{漢}y$漢
\ifnum\jlreqpdf@math@rawerrcnt=1
  \typeout{TATE-PASS: RAWERR (raw hbox raised ClassError)}%
\else
  \typeout{TATE-FAIL: RAWERR (count=\the\jlreqpdf@math@rawerrcnt, expected 1)}%
\fi
% ---- VBOXERR: raw \vbox → ClassError ----
漢$x\vbox{漢}y$漢
\ifnum\jlreqpdf@math@rawerrcnt=2
  \typeout{TATE-PASS: VBOXERR (raw vbox raised ClassError; total 2)}%
\else
  \typeout{TATE-FAIL: VBOXERR (count=\the\jlreqpdf@math@rawerrcnt, expected 2)}%
\fi
\end{document}
